\documentclass[12pt]{article}
\usepackage{amsmath}
\usepackage{enumitem}
\usepackage{geometry}
\geometry{margin=1in}

\title{Planuri de Execuție în SQL}
\author{}
\date{}

\begin{document}
\maketitle

\section{Ce este un Plan de Execuție}
Un plan de execuție reprezintă blueprint-ul intern al motorului SQL care descrie:
\begin{itemize}[noitemsep]
    \item modul de accesare a datelor (Index Scan, Full Table Scan)
    \item strategia de îmbinare a tabelelor (Nested Loop, Hash Join, Merge Join)
    \item ordinea operațiilor
    \item costul estimat al fiecărei operații
    \item cardinalitatea estimată (numărul de rânduri procesate)
\end{itemize}

Planul este generat de optimizatorul bazei de date şi poate fi afişat prin comenzi precum \texttt{EXPLAIN}, \texttt{EXPLAIN ANALYZE} sau \texttt{DBMS\_XPLAN.DISPLAY}.

\section{Tipuri de Operații în Planurile de Execuție}

\subsection{Full Table Scan}
\textbf{Descriere:} Motorul citeşte întregul tabel rând cu rând.  
\textbf{Utilizare:} când nu există index util sau filtrul nu este selectiv.

\textbf{Pseudo-implementare:}
\begin{verbatim}
for row in Table:
    if condition(row):
        output(row)
\end{verbatim}

\subsection{Index Scan}
\textbf{Descriere:} Parcurge indexul secvenţial, nu direct pe cheie.  
\textbf{Utilizare:} când filtrul foloseşte o coloană indexată, dar intervalul este larg.

\textbf{Pseudo-implementare:}
\begin{verbatim}
for entry in Index:
    if condition(entry.key):
        fetch row(entry)
\end{verbatim}

\subsection{Index Seek}
\textbf{Descriere:} Acces direct în index prin căutare binară.  
\textbf{Utilizare:} cel mai rapid acces, când filtrul este foarte selectiv.

\textbf{Pseudo-implementare:}
\begin{verbatim}
node = BTree.root
while node not leaf:
    node = choose_child(node, search_key)
return fetch_rows(node, search_key)
\end{verbatim}

\subsection{Nested Loop Join}
\textbf{Descriere:} Pentru fiecare rând din tabelul A, caută rândurile corespunzătoare în tabelul B.  
\textbf{Utilizare:} seturi mici sau când B are index pe coloana de join.

\textbf{Pseudo-implementare:}
\begin{verbatim}
for rowA in A:
    for rowB in B where rowB.key = rowA.key:
        output(rowA, rowB)
\end{verbatim}

\subsection{Hash Join}
\textbf{Descriere:} Construieşte un hash table pentru tabelul mai mic, apoi îl sondează cu rândurile din tabelul mare.  
\textbf{Utilizare:} seturi mari, fără indexuri.

\textbf{Pseudo-implementare:}
\begin{verbatim}
hash = {}
for rowA in A:
    hash[rowA.key].append(rowA)

for rowB in B:
    if rowB.key in hash:
        for rowA in hash[rowB.key]:
            output(rowA, rowB)
\end{verbatim}

\subsection{Merge Join}
\textbf{Descriere:} Ambele tabele sunt sortate după cheia de join, apoi parcurse simultan.  
\textbf{Utilizare:} când datele sunt deja sortate sau există indexuri ordonate.

\textbf{Pseudo-implementare:}
\begin{verbatim}
i = j = 0
while i < len(A) and j < len(B):
    if A[i].key == B[j].key:
        output(A[i], B[j])
        i += 1; j += 1
    elif A[i].key < B[j].key:
        i += 1
    else:
        j += 1
\end{verbatim}

\subsection{Sort}
\textbf{Descriere:} Operaţie costisitoare, necesară pentru ORDER BY, GROUP BY sau Merge Join.

\textbf{Pseudo-implementare:}
\begin{verbatim}
sorted_rows = quicksort(rows)
\end{verbatim}

\subsection{Filter}
\textbf{Descriere:} Aplică predicatul din clauza WHERE.

\textbf{Pseudo-implementare:}
\begin{verbatim}
for row in rows:
    if condition(row):
        output(row)
\end{verbatim}

\section{Hinturi SQL}

Hinturile SQL sunt instrucțiuni speciale inserate în comentarii, folosite pentru a influenţa modul în care optimizatorul generează planul de execuţie. Ele nu modifică logica interogării, ci doar strategia de execuţie.

\subsection{Oracle SQL}

Hinturile se inserează imediat după cuvântul \texttt{SELECT} sub forma unui comentariu special:

\begin{verbatim}
SELECT /*+ HINT_NAME */ coloane
FROM tabel;
\end{verbatim}

Exemple:

\begin{itemize}
    \item \textbf{Forţare Index}
\begin{verbatim}
SELECT /*+ INDEX(emp emp_idx) */ *
FROM employees emp
WHERE emp.department_id = 50;
\end{verbatim}

    \item \textbf{Forţare Full Table Scan}
\begin{verbatim}
SELECT /*+ FULL(emp) */ *
FROM employees emp;
\end{verbatim}

    \item \textbf{Forţare Hash Join}
\begin{verbatim}
SELECT /*+ USE_HASH(e d) */ e.name, d.name
FROM employees e
JOIN departments d ON e.dept_id = d.id;
\end{verbatim}

    \item \textbf{Forţare Nested Loop}
\begin{verbatim}
SELECT /*+ USE_NL(e d) */ e.name, d.name
FROM employees e
JOIN departments d ON e.dept_id = d.id;
\end{verbatim}
\end{itemize}

\subsection{SQL Server}

SQL Server nu foloseşte hinturi în comentarii, ci direct în sintaxă:

\begin{verbatim}
SELECT coloane
FROM tabel WITH (INDEX(index_name));
\end{verbatim}

Exemple:

\begin{itemize}
    \item \textbf{Forţare Index}
\begin{verbatim}
SELECT *
FROM Employees WITH (INDEX(emp_idx))
WHERE DepartmentID = 50;
\end{verbatim}

    \item \textbf{Forţare NOLOCK}\\
Hintul \texttt{NOLOCK} permite citirea datelor fără a achiziţiona shared locks, ignorând astfel blocările existente. Acest lucru poate îmbunătăţi concurenţa şi reduce blocajele, însă poate produce citiri incorecte precum dirty reads, non-repeatable reads sau phantom reads.
\begin{verbatim}
SELECT *
FROM Employees WITH (NOLOCK);
\end{verbatim}
\end{itemize}

\subsection{PostgreSQL}

PostgreSQL nu suportă hinturi nativ.  
Hinturile pot fi folosite doar prin extensia \texttt{pg\_hint\_plan}:

\begin{verbatim}
SELECT /*+ IndexScan(employees emp_idx) */ *
FROM employees
WHERE department_id = 50;
\end{verbatim}


\end{document}
